\section{Discussion}

\subsection{Key Findings}

Our experiments reveal three important findings about cost-performance trade-offs in UQ for LLM selective prediction:

\textbf{Finding 1: Epistemic uncertainty threshold exists at 8B scale.} MC dropout $k\geq3$ required for AUROC $\geq$ 0.70, while zero-cost methods (temperature scaling, conformal prediction) fall below threshold. This suggests post-hoc calibration (aleatoric uncertainty) is insufficient when base model has complex miscalibration patterns (TruthfulQA adversarial questions). Epistemic uncertainty---captured via stochastic forward passes---provides stronger signal for rejecting incorrect predictions.

This finding challenges the assumption that calibration methods are interchangeable. At 8B scale, \textbf{the mechanism matters}: temperature scaling rescales logits assuming monotonic confidence-correctness relationship, but TruthfulQA adversarial design breaks this assumption. MC dropout samples from approximate Bayesian posterior, directly quantifying regions where model is uncertain. The 3\% AUROC gap (0.682 vs 0.712) represents the value of epistemic uncertainty for high-stakes selective prediction.

\textbf{Finding 2: MC dropout $k=3$ is an efficiency sweet spot.} $k=3$ achieves 0.704 AUROC at 3$\times$ cost, offering 40\% cost savings vs $k=5$ (0.712 AUROC at 5$\times$ cost) for only 0.008 AUROC drop. This contrasts with vision domain conventions ($k=30$ in retinal-selective-prediction) and suggests Bayesian approximation converges faster at 8B model scale. Smaller parameter space $\rightarrow$ faster posterior convergence, making high $k$ unnecessary.

For practitioners, this means \textbf{$k=3$ should be the production default} for applications requiring $\geq$0.70 AUROC. The marginal AUROC gain from $k=5$ (+0.008) rarely justifies doubling cost from 3$\times$ to 5$\times$, unless application is truly high-stakes (e.g., medical diagnosis, legal advice).

\textbf{Finding 3: Pareto frontier validates budget-aware UQ selection.} 5 methods occupy distinct positions---no universal dominance across budget constraints. This shifts the question from ``which method is best?'' to ``which method for my budget?'', opening research directions in adaptive $k$ methods (early stopping when variance converges), hybrid approaches (cheap method for easy queries, expensive for hard queries), and per-query budget allocation.

The Pareto framing also reveals when \textbf{zero-cost methods are acceptable}: applications with <3$\times$ budget constraint AND tolerance for <0.70 AUROC (e.g., content moderation with human review fallback) can use temperature scaling or conformal prediction. The key is making this trade-off explicit rather than defaulting to expensive methods everywhere.

\subsection{Limitations}

Our work has several limitations:

\textbf{Limitation 1: 8B model scale only.} Results are specific to Llama-3.1-8B-Instruct. Larger models (70B, 405B) may have better base calibration, potentially enabling zero-cost methods to exceed 0.70 threshold. Conversely, smaller models (GPT-2, 1B) may require higher $k$ for threshold crossing.

\textbf{Why acceptable:} 8B scale is realistic for many production deployments (latency/memory constraints). Our findings establish feasibility at this scale and identify $k=3$ sweet spot---useful even if 70B results differ.

\textbf{Future work:} Repeat h-e1/h-m-pareto on Llama-3.1-70B-Instruct to test if zero-cost methods become competitive at larger scale.

\textbf{Limitation 2: Single benchmark (TruthfulQA).} Task-dependent AUROC thresholds exist---MMLU (knowledge QA) may be easier, GSM8K (math reasoning) harder. The 0.70 threshold may not generalize across domains.

\textbf{Why acceptable:} TruthfulQA adversarial design provides challenging testbed (31\% base accuracy). Methods working here likely transfer to easier tasks. Threshold can be task-adjusted (e.g., 0.65 for GSM8K, 0.75 for MMLU).

\textbf{Future work:} Cross-benchmark validation (MMLU, GSM8K, HaluEval) to characterize task-dependent thresholds and method ranking stability.

\textbf{Limitation 3: $n=1$ seed for full-scale validation (low statistical power).} MC $k=10$ vs $k=5$ difference (+0.006 AUROC) not statistically significant with single seed. h-m-pareto PoC used $n=3$ seeds but on smaller GPT-2 model.

\textbf{Why acceptable:} Small standard deviation ($\pm$0.0082 from PoC) suggests difference is genuine though underpowered. Increasing to $n=5$ seeds would cost 5$\times$ compute for marginal confidence gain.

\textbf{Future work:} $n=5$ seeds for 70B model study to achieve 80\% statistical power for pairwise comparisons.

\textbf{Limitation 4: HaluEval $\rightarrow$ TruthfulQA calibration transfer gap.} Conformal prediction AUROC 0.695 (-0.005 below threshold) due to domain shift (hallucination detection $\neq$ truthfulness QA). In-distribution calibration expected to improve AUROC.

\textbf{Why acceptable:} Cross-dataset transfer tests generalization, which is valuable for practitioners. Coverage guarantees still hold per conformal theory---AUROC is quality metric, not validity metric.

\textbf{Future work:} Use TruthfulQA 40\% calibration split instead of HaluEval to close gap and test if conformal reaches threshold.

\subsection{Broader Impact}

\textbf{Positive impacts:} Budget-aware UQ selection enables more practitioners to deploy selective prediction in production (not just research labs with unlimited compute). $k=3$ sweet spot reduces barrier to entry for resource-constrained teams.

\textbf{Negative impacts:} Practitioners may over-optimize for cost at expense of safety. Applications requiring high confidence (medical, legal) should not use $k=3$ if $k=5$ is affordable---our framework guides trade-offs but does not mandate specific choices.

\textbf{Mitigation:} Clear communication that 0.70 threshold is baseline for viable selective prediction, not guarantee of safety. High-stakes applications should use $k\geq5$ and validate on task-specific benchmarks.

\subsection{Paradigm Shift Potential}

If replicated at 70B scale showing zero-cost methods exceed 0.70 threshold, this would reverse the ``more compute = better UQ'' narrative. Practitioners could \textbf{default to temperature scaling} (0$\times$ overhead) unless high-stakes application requires MC dropout. This shift could reduce aggregate compute waste in production LLM deployments.

Conversely, if 70B results confirm epistemic uncertainty requirement, the field should standardize on \textbf{$k=3$ as production default} rather than $k=30$ conventions from vision domain. Our Pareto frontier provides the cost-performance map needed for this standardization.
